The ALICE Event Processing Nodes, the farm machines that reconstruct the experiment's data, write logs that mostly never reach the tool people search. The goal of this project was a log platform that keeps those logs cheap for the workers and still lets one person search and watch the whole farm. \loggy{} does this. A collector on each worker stamps every line with its message template. Info, debug and trace lines stay on the worker that wrote them. Warnings, errors and InfoLogger records go to three storage machines. All machines form one OpenSearch cluster, so one query searches every worker. Each worker pushes its health up, and threshold rules, anomaly detectors and trend rules turn those numbers into episodes a person is told about. \loggy{} runs on five staging machines, and a farm deployment of the cluster, the collectors and Dashboards passed every check. In six months of archive, the busiest second any worker had carried 78~records. The whole stack on one worker sustains about 42,000 records a second. Every cost figure comes from a laptop, not from a farm machine. The shape of each curve and the ranking of settings transfer to the farm. The absolute rates need one benchmark on a farm machine, and that benchmark has not run.
